\section{Tái lập}
\label{app:repro}

\subsection{Môi trường và phần cứng}

\begin{itemize}
    \item Python 3.13.5. Thư viện (\texttt{requirements.txt}): numpy 2.3.1, pandas 2.3.1, scipy 1.18.1, matplotlib 3.11.0, scikit-learn 1.9.0.
    \item PyTorch 2.14.0. Bản \texttt{+cu130} cho GPU RTX 50xx (Blackwell cần CUDA $\ge$ 12.8); bản CPU nếu không có GPU. Cài torch \emph{trước} \texttt{requirements.txt}, xem \texttt{README.md}.
    \item GPU (grid, IRLS, code gốc LTSF-Linear): NVIDIA RTX 5060 Ti 16 GB, driver 616.92, CUDA 13.0, TF32 tắt. CPU (SGD): Intel 12 luồng. Hệ điều hành Windows 11; lệnh chạy qua Git Bash.
\end{itemize}

\subsection{Seed}

\begin{itemize}
    \item SGD với lr gốc: seed 2021--2040 (NLinear trên ETTh2: 2021--2030).
    \item SGD tinh chỉnh lr: chọn lr bằng seed 2021--2025, chạy cuối bằng seed 2026--2040 (không trùng seed chọn lr).
    \item Code gốc LTSF-Linear (có và không \texttt{drop\_last}): seed 2021--2025.
    \item Vị trí outlier: seed 2026; thêm 2027, 2028 ở ETTh1 $H = 96$ Linear.
\end{itemize}

\subsection{Lệnh}

\begin{aivncodebox}[Tái lập]
\begin{python}
git clone --recurse-submodules <repo> && cd dlinear-objective-lab
./run_all.sh quick   # tests, re-run one grid cell, regenerate numbers/tables/figures, build PDF
./run_all.sh full    # re-run every experiment (many GPU/CPU hours)
\end{python}
\end{aivncodebox}

Chế độ \texttt{quick} chạy lại toàn bộ grid của ETTh1 $H = 96$ (3 mô hình, 3 mục tiêu) vào một thư mục tạm, rồi so với \texttt{results/lambda\_path.csv} bằng \texttt{scripts/checks/compare\_cell.py}. Ngưỡng là $10^{-10}$ với MSE và $10^{-8}$ với MAE, Huber. Trên máy ở trên, bước này mất khoảng \nmReproCheckCellMinutes{} phút. Chế độ \texttt{full} chạy theo thứ tự: grid, $\lambda^*$, Huber $\dmad$, outlier, SGD, tinh chỉnh lr và mở rộng lưới lr, code gốc với \texttt{drop\_last}, Huber với $\delta$ theo bước. Thời gian từng bước, giá trị mong đợi của ô kiểm và lệnh chi tiết nằm trong \texttt{README.md} và mục ``Tái lập'' của từng báo cáo kỹ thuật ở \texttt{docs/reports/}. Đường dẫn cũ trong các báo cáo đó tra ở \texttt{docs/DUONG\_DAN\_CU.md}.

\subsection{Thời gian chạy đã đo}

\begin{itemize}
    \item Grid (MSE dạng đóng, MAE và Huber $\delta = 1$ bằng IRLS, GPU): \nmReproGridHours{} giờ, gồm ETTh1 \nmReproGridHoursETThOne, ETTh2 \nmReproGridHoursETThTwo{} và ETTm1 \nmReproGridHoursETTmOne{} giờ. MAE chiếm gần hết thời gian.
    \item Grid Huber $\dmad$ (GPU): \nmReproHuberMadHours{} giờ.
    \item SGD với lr gốc (CPU, tổng thời gian của các lần chạy, 20 seed): \nmReproSgdCpuHours{} giờ.
    \item Tinh chỉnh lr (CPU, 5 tiến trình song song): \nmReproSgdLrWallHours{} giờ thời gian thực; mở rộng lưới ở các ô đầu mút: \nmTunedLrExtWallHours{} giờ.
    \item Code gốc LTSF-Linear với \texttt{drop\_last} bật và tắt (GPU, tổng thời gian của các lần chạy): \nmReproDropGpuHours{} giờ.
    \item Huber với $\delta$ theo bước, 2 ô (GPU): \nmReproDeltaHMinutes{} phút.
\end{itemize}

\subsection{Phiên bản}

Các bảng, hình và số trong báo cáo được sinh từ \texttt{results/} ở commit \nmReproHead{} (commit cuối cùng thay đổi \texttt{results/}). Mỗi dòng trong các file \texttt{lambda\_path.csv} và mỗi JSON của SGD ghi commit đã sinh ra nó (cột \texttt{git\_commit}). LTSF-Linear được ghim ở commit \texttt{0c11366} qua git submodule.
